\section{Matrices}

A matrix is a rectangular array of numbers. For example,
\[
A=\begin{pmatrix}
1 & 2 \\
3 & 4
\end{pmatrix}.
\]

\begin{center}
\begin{tabular}{@{}ll@{}}
\toprule
Operation & Example \\
\midrule
Addition & \(\begin{pmatrix}1&2\\3&4\end{pmatrix}+\begin{pmatrix}5&6\\7&8\end{pmatrix}=\begin{pmatrix}6&8\\10&12\end{pmatrix}\) \\
Scalar multiplication & \(2\begin{pmatrix}1&2\\3&4\end{pmatrix}=\begin{pmatrix}2&4\\6&8\end{pmatrix}\) \\
Matrix times vector & \(\begin{pmatrix}1&2\\3&4\end{pmatrix}\begin{pmatrix}x\\y\end{pmatrix}=\begin{pmatrix}x+2y\\3x+4y\end{pmatrix}\) \\
\bottomrule
\end{tabular}
\end{center}

Matrix multiplication is not the same as ordinary multiplication of numbers. In general,
\[
AB \neq BA.
\]
This means that the order of multiplication matters.

Matrices are useful because they can represent linear transformations. For example, the matrix
\[
R(\theta)=
\begin{pmatrix}
\cos\theta & -\sin\theta \\
\sin\theta & \cos\theta
\end{pmatrix}
\]
represents a rotation in the plane by an angle \(\theta\). If
\[
\mathbf{x}=\begin{pmatrix}x\\y\end{pmatrix},
\]
then \(R(\theta)\mathbf{x}\) is the rotated vector.

\begin{remarkbox}
In these notes, matrices will not be developed in detail. The reader should remember the basic idea: matrices organize several numbers and can describe transformations acting on vectors or systems of quantities.
\end{remarkbox}
